\documentclass[11pt,letterpaper]{article}
\usepackage[T1]{fontenc}
\usepackage{lmodern,amsmath,amssymb,amsthm,microtype}
\usepackage[margin=1in]{geometry}
\usepackage[hidelinks]{hyperref}
\providecommand{\tightlist}{\setlength{\itemsep}{0pt}\setlength{\parskip}{0pt}}
\title{The four-isolate case of the height-six Frankl model}
\author{Aafi Khan Malek\\ \small Independent researcher, Ahmedabad, India\\
\small\href{mailto:aafimalek2032@gmail.com}{aafimalek2032@gmail.com}}
\date{28 November 2026}
\DeclareUnicodeCharacter{222A}{\ensuremath{\cup}}
\DeclareUnicodeCharacter{2208}{\ensuremath{\in}}
\begin{document}
\maketitle
\textbf{Written argument (not checked end to end in Lean).} Under the
same weak frequency, minimum-ground, and height-six hypotheses as the
four-isolates bound, a family with exactly four isolates must have
exactly one forcing source. A separate finite trace calculation below
reduces that remaining case to one explicit trace family, but does not
rule it out.

Throughout, \texttt{k} counts sources, \texttt{N} counts blocks of the
omitted-star pair design, \texttt{Q=binom(k,2)}, and
\texttt{H(m)=1+m+binom(m,2)}. A block of size \texttt{r} incurs
lost-block count \texttt{binom(r,2)-1} from the all-pair design. Write
\texttt{d\_v=k-1-g\_v} for the deficit of vertex \texttt{v}. Pair
uniqueness implies

\begin{verbatim}
 d_v = sum_{B contains v} (|B|-2).
\end{verbatim}

If a pair-avoidance fiber is supported, after forced identifications, on
at most \texttt{m} coordinates, then its cardinality is at most
\texttt{H(m)}. Its existing witness lower bound (using ground size at
least \texttt{2k+4}) for anchors \texttt{p,q} is

\begin{verbatim}
 |F(p,q)| >= N+k+6+d_p+d_q.                           (A)
\end{verbatim}

This is \texttt{n+k+(N-g\_p-g\_q+1)-1} with
\texttt{n\textgreater{}=2k+4}.

Zero sources give a four-point ground and are impossible by the external
12-point theorem, as are all ground sizes at most twelve: adjoin the
empty set to convert a no-strict-majority family into an ordinary
below-half counterexample. If there are two or more sources and a
nonunary star, the earlier multileaf lemma gives
\texttt{n+k\textless{}=12} at four isolates, again impossible. Hence in
the rest of the multiple-source proof all stars are unary,
\texttt{n=2k+4}, and \texttt{k\textgreater{}=5}.

\section{Maximum degree deficit at most
two}\label{1-maximum-degree-deficit-at-most-two}

Choose a minimum degree \texttt{g} at \texttt{p}, put \texttt{d=k-1-g},
and anchor with any distinct \texttt{q}. The checked block link groups
from \texttt{p} give
\texttt{\textbar{}F(p,q)\textbar{}\textless{}=H(g+3)}. From (A) the
lower bound is at least \texttt{N+2k+5-g}. Substitute the earlier
incidence bound \texttt{N\textgreater{}=k*g\^{}2/(k-1+g)}. Clearing the
positive denominator yields

\begin{verbatim}
 (d-8)g^2+3dg+4d^2 <= 0.
\end{verbatim}

Since \texttt{g\textgreater{}=1}, this implies \texttt{d\textless{}=7}.
The block contribution \texttt{binom(r,2)-1\ \textless{}=\ 2r(r-2)/3}
and \texttt{sum\_v\ d\_v\textless{}=kd} give
\texttt{N\textgreater{}=Q-2kd/3}. Substituting this instead bounds the
available slack between the upper and lower bounds by

\begin{verbatim}
 (2-d/3)k+(d^2-7d-4)/2.
\end{verbatim}

This is negative at \texttt{d=6,7}. For
\texttt{3\textless{}=d\textless{}=5}, even one merger of the available
anchor groups costs \texttt{H(g+3)-H(g+2)=g+3}, which exceeds that
slack: the differences for \texttt{d=3,4,5} are respectively \texttt{7},
\texttt{k/3+6}, and \texttt{2k/3+4}. Thus every exterior group through
\texttt{q} is contained in an exterior group through \texttt{p}. Pair
uniqueness makes every nonpair block incident to \texttt{p}, exactly as
in the three-isolate argument. For clarity: if a non-anchor block
through \texttt{q} had two exterior points, they would also belong to
one different block through \texttt{p}, repeating their pair. The
conclusion holds for arbitrary \texttt{q\ !=\ p}.

With all nonpair blocks through \texttt{p}, their positive contributions
\texttt{e=\textbar{}B\textbar{}-2} sum to \texttt{d}. Their lost-block
count is \texttt{(e\^{}2+3e)/2}; convexity (or
\texttt{sum\ e\^{}2\textless{}=d\^{}2}) gives
\texttt{N\textgreater{}=Q-(d\^{}2+3d)/2}. The available slack is then at
most \texttt{(d-2)(d-k)-2\textless{}0} for
\texttt{3\textless{}=d\textless{}=5}. Therefore \texttt{d\textless{}=2}.

\section{A reusable near-extremal counting
observation}\label{2-a-reusable-near-extremal-counting-observation}

After identifying the forced equal coordinates in a pair fiber,
complement its members relative to its true support. This gives an
intersection-closed family with the empty set and no four-member strict
chain. The rank-two atom bound proved in the preceding note gives, on at
most \texttt{m} coordinates, at most \texttt{1+m+binom(a,2)} members,
where \texttt{a} is the number of disjoint nonempty atoms. Thus if the
deficit from \texttt{H(m)} is less than \texttt{m-1}, all \texttt{m}
atoms are singletons. In particular all coordinates are genuinely
available, the fiber maximum is full support, and all relative singleton
omissions belong to the complemented fiber. If only
\texttt{a\textgreater{}=m-1} is known and the true support is all
\texttt{m} coordinates, all but at most two coordinates are singleton
atoms; if two fail to be singleton atoms they form a two-coordinate
atom. If one available coordinate is absent from the support, the true
support has exactly \texttt{m-1} coordinates and all its atoms are
singletons.

Call a ground complement \texttt{S\_r∪J} a mixed omission if it is the
complement of some family member, where \texttt{S\_r} is a unary star
and \texttt{J} is a nonempty isolate subset. When \texttt{r} is outside
the anchor pair, this counts as an extra witness in (A): the member
contains both anchor leaves, but is distinct from the ground,
single-point avoidance maxima (at least three original points are
omitted), source witness maxima (omits a leaf), and pair maxima (omits
only one whole star). If two different pair blocks \texttt{\{r,a\}} and
\texttt{\{r,b\}} have full available support and both possess the
\emph{same} relative isolate omission \texttt{J}, unioning those two
members gives the mixed omission \texttt{S\_r∪J}.

\section{Deficit zero: all-pair
design}\label{3-deficit-zero-all-pair-design}

Here every anchor pair has \texttt{m=k+2} available coordinates before
optional identifications. Its bound from below is \texttt{Q+k+6}, while
\texttt{H(k+2)=Q+3k+4}, leaving slack \texttt{2k-2}. If
\texttt{a\textless{}=k}, the atom bound gives at most
\texttt{1+(k+2)+Q=Q+k+3}, impossible. Thus every pair has at least
\texttt{m-1} atoms and at least two of the four isolates are singleton
atoms.

For each source \texttt{r}, consider the \texttt{k-1} anchor pairs
through \texttt{r}. For a fixed such pair the exact original isolate
omissions, in members omitting just those two stars, include one of the
following three lists:

\begin{enumerate}
\def\labelenumi{\arabic{enumi}.}
\tightlist
\item
  At least three singleton isolates (full support).
\item
  Exactly two singleton isolates and their complementary two-isolate
  set, which is the two-coordinate atom (full support).
\item
  Some singleton \texttt{\{z\}} and all pairs \texttt{\{z,w\}} with
  \texttt{w!=z}: this is the case of a missing coordinate. It must be an
  isolate \texttt{z}, since the maximum omits exactly the anchor stars.
  All other coordinates are singleton atoms on the true support, giving
  the displayed omissions by deleting those coordinates or none.
\end{enumerate}

The union of two such members with different anchor partners has
complement the star \texttt{S\_r} together with the intersection of
their isolate omissions. We claim at least three distinct nonempty
intersections for each \texttt{r}.

\begin{itemize}
\tightlist
\item
  If two lists are of type 3 with the same missing isolate, their common
  list already gives four. With different missing isolates \texttt{z,w},
  the cross-intersections give \texttt{\{z\},\{w\},\{z,w\}}.
\item
  If exactly one list is of type 3, every singleton isolate omission of
  another list can be intersected with a type 3 omission to obtain that
  singleton. If only two singletons occur in all other lists, all those
  lists have type 2 with a common complementary pair \texttt{J}.
  Intersect \texttt{J} with a suitable type 3 omission to obtain any one
  singleton in \texttt{J}.
\item
  If none is of type 3, at least two isolates appear as singletons in at
  least two lists apiece: otherwise there are at most
  \texttt{(k-1)+3\ \textless{}\ 2(k-1)} singleton occurrences. If three
  or more are repeated, their singleton intersections suffice. If only
  two are repeated, the other two appear a combined total at most twice.
  At least \texttt{k-3\textgreater{}=2} lists are then type 2 with the
  same complementary pair of these two non-repeated isolates. This
  yields the third nonempty intersection.
\end{itemize}

Therefore at least three mixed omissions are available for each outside
star. For any fixed anchor pair the resulting
\texttt{3(k-2)\textgreater{}2k-2} extra witnesses exceed its slack, a
contradiction.

\section{Deficit one: disjoint
triples}\label{4-deficit-one-disjoint-triples}

If the maximum deficit is one, the only nonpair blocks are vertex-
disjoint triples. Let there be \texttt{t\textgreater{}=1}, so
\texttt{N=Q-2t}. For anchors in two distinct triples the anchor block is
a pair, and the two disjoint mate pairs each collapse to one coordinate.
Thus \texttt{m=k}. Their lower bound is \texttt{Q-2t+k+8}, so the slack
in \texttt{H(k)=Q+k+1} is \texttt{2t-7}. This is negative for
\texttt{t=2,3}; when \texttt{t\textgreater{}=4} it is less than
\texttt{k-1}, since \texttt{3t\textless{}=k}. Thus all quotient
singleton omissions exist and the maximum has full support on the
quotient. For any vertex in a triple and any isolate, taking two
partners from another triple gives the mixed omission of its star and
that isolate. This provides \texttt{4(3t-2)\textgreater{}2t-7} extra
witnesses for a fixed distinct-triple anchor pair, impossible.

If \texttt{t=1}, use an anchor consisting of a vertex of that triple and
one outside it. The anchor block is a pair and the only required
identification is one mate pair, so \texttt{m=k+1}. The lower bound is
\texttt{Q+k+5}; slack \texttt{k-3} is less than \texttt{m-1}. Again all
quotient singletons exist. Every source can take two partners from the
opposite side of the partition into triple and outside (the outside has
at least two vertices because \texttt{k\textgreater{}=5}). It follows
that every source gives all four mixed singleton isolate omissions. The
resulting \texttt{4(k-2)\textgreater{}k-3} witnesses exceed the slack.

\section{Deficit two: triples and
four-blocks}\label{5-deficit-two-triples-and-four-blocks}

Now each four-block is disjoint from all other nonpair blocks: each of
its vertices is already at deficit two. A vertex outside the four-blocks
belongs to at most two triples. If there are \texttt{t} triples and
\texttt{u} four-blocks, put \texttt{S=Q-N=2t+5u}. Counting total
deficits gives \texttt{3t+8u\textless{}=2k}, hence
\texttt{S\textless{}=4k/3}.

Two distinct deficit-two anchors have lower bound \texttt{Q-S+k+10},
which is greater than \texttt{H(k-2)=Q-k+2}. Thus any geometry
compressing their fiber to at most \texttt{k-2} coordinates is
impossible.

For such anchors in different blocks (their anchor block is a pair),
start with \texttt{k+2} available coordinates. Each anchor contributes
two units of identification among outsider stars: either one three-point
group (for a four-block), or two disjoint two-point groups (for two
triples). Groups from opposite anchors meet in at most one point, by
pair uniqueness. The four units of identification are independent unless
both anchors consist of two mate pairs and the four pairs make a cycle
of length four. To see this, form a bipartite graph with the two
anchors\textquotesingle{} groups as vertices and an edge for each star
in both groups. Groups on each side are disjoint, and the graph is
simple by pair uniqueness. The number of identifications contributed
independently is four minus its cycle rank: adding an identification
within a group joins two coordinate classes unless they were already
joined via a sequence of groups on the opposite side. Equivalently, in
each connected component with \texttt{v} groups and \texttt{e} pairwise
overlaps, the number of distinct covered coordinates is
\texttt{sum(group\ sizes)-e}; the forced reduction there is that number
minus one, less than the sum of the group reductions by \texttt{e-v+1}.
With at most two groups on each side the only possible cycle is the
four-cycle, and it requires both sides to have two groups. Without it
there are four identifications, which is impossible by the preceding
cardinality observation.

\subsection{A four-block exists}\label{a-four-block-exists}

No deficit-two vertex can lie outside the four-block, since the
preceding argument would require its three-point group to be on a
four-cycle. Thus there is one four-block, and all other nonpair blocks
are disjoint triples. If any triple exists, take one anchor in the
four-block and one in the triple. Their group reductions are independent
and total three, so \texttt{m\textless{}=k-1}. But
\texttt{S=5+2t\textless{}=5+2(k-4)/3\textless{}k+8}, and (A) gives
\texttt{Q-S+k+9\textgreater{}Q+1=H(k-1)}, impossible. If no triple
exists, two anchors inside the four-block have \texttt{m=k} and (A)
gives \texttt{Q-5+k+10\textgreater{}H(k)}, also impossible.

\subsection{Only triples, with at least one deficit-two
source}\label{only-triples-with-at-least-one-deficit-two-source}

Make a simple graph with one vertex for each triple and one edge for
each source belonging to two triples, joining those two triples. It has
maximum degree three, since a triple has three sources. Any two disjoint
edges correspond to separate deficit- two anchors whose mate pairs must
form the above four-cycle. Each of the four intersections gives a source
incident with one triple from each edge, hence the four cross-edges of a
literal \texttt{K4}. Its four vertices are cubic within the \texttt{K4},
so any further edge would be disjoint from it. Such an edge, together
with an edge from the \texttt{K4}, would require cross-edges, which is
impossible. If there are no disjoint edges, the family of edges is
pairwise intersecting; in a simple graph this is a star or a triangle.
With maximum degree three it has at most three edges. In either case the
number \texttt{D} of deficit-two sources is at most six.

If there are at least two deficit-two sources, take any two. If they
belong to a common triple, that block removes three stars, leaving
\texttt{k+1} available coordinates; their two different outer mate pairs
give two independent identifications. If they do not, their mates give
the four-cycle and three independent identifications on \texttt{k+2}
coordinates. Thus \texttt{m\textless{}=k-1} in both cases. From (A),
\texttt{S\textgreater{}=k+9} is necessary, but
\texttt{3t\textless{}=k+D\textless{}=k+6} gives
\texttt{S=2t\textless{}=2(k+6)/3\textless{}k+9}, a contradiction.

Finally suppose there is exactly one deficit-two source \texttt{p}. If a
triple other than its two exists, it is disjoint from both; a source
\texttt{q} of that triple has deficit one. Their three independent group
identifications give \texttt{m\textless{}=k-1}; (A) would require
\texttt{S\textgreater{}=k+8}, whereas \texttt{3t\textless{}=k+1} implies
\texttt{S=2t\textless{}=2(k+1)/3}. If the only triples are the two
through \texttt{p}, choose a mate of \texttt{p} in one triple as the
other anchor. Removing the anchor triple leaves \texttt{k+1}
coordinates, and the other mate pair collapses once, so \texttt{m=k}.
But (A) gives \texttt{Q-4+k+9\textgreater{}H(k)}.

All three possible maximum deficits are excluded, completing the
multiple-source claim.

\section{Remaining one-source trace family: an exact finite
reduction}\label{6-remaining-one-source-trace-family-an-exact-finite-reduction}

Let \texttt{I} be the four isolates and \texttt{L} count the leaves at
the unique source. Let \texttt{P} be the family of traces on \texttt{I}
of the source-containing members, and let \texttt{B} be the
source-avoiding family. Every source-containing set contains the entire
star; the checked top and co-singleton facts put all 3- and 4-subsets of
\texttt{I} in \texttt{P}. The checked smaller-ground reduction gives an
isolate \texttt{x} of strictly below-half frequency in \texttt{P};
rename it coordinate zero. The trace of any member of \texttt{B} must be
one of the safe traces \texttt{T}, defined by \texttt{T∪R\ in\ P} for
every \texttt{R∈P}.

The self-contained exact replay \texttt{traces.py} exhausts the
\texttt{2\^{}11} choices for the optional traces, rejecting families not
union-closed or without the requisite low point. There are 106 remaining
labeled families. For each candidate and each pair of isolates
\texttt{y,z}, it computes

\begin{verbatim}
 s = |union of all safe traces avoiding y,z|;
 a = frequency_P(y)+frequency_P(z)-|P|.
\end{verbatim}

If \texttt{a\textgreater{}s} and the safe supporting union is contained
in the bottom trace of a four-member chain in \texttt{P}, the candidate
is impossible. Indeed the source-avoiding family omitting both
\texttt{y,z} can extend any chain by those four source-containing
members, and thus has height at most two on at most \texttt{L+s} points.
Its cardinality is at most \texttt{L+s+1}. Meanwhile there are at least
\texttt{L+1} source- avoiding members containing both isolates (the full
source- avoiding top and its co-leaf tops). Inclusion--exclusion and the
two weak-half bounds give at least \texttt{L+1+a} members omitting both,
a contradiction.

The replay excludes 105 and leaves exactly the trace family
\texttt{P=2\^{}I\ \textbackslash{}\ \{\{x\}\}}. Its four frequencies are
\texttt{7,8,8,8} out of 15. This model is not excluded by the
calculation: no claimed contradiction for the exactly-four-isolate,
one-source case.

For this residual, write \texttt{b=\textbar{}B\textbar{}}. The high
isolates each occur at most \texttt{(b-1)/2} times in \texttt{B}, and
the low isolate \texttt{x} is the only possible strict majority in
\texttt{B} by smaller ground minimality (every leaf occurs in all 15
source-containing sets). The full-family bound then gives \texttt{b} odd
with \texttt{f\_B(x)=(b+1)/2}. Trace \texttt{\{x\}} is forbidden. The
empty-trace subfamily of \texttt{B} has at most one member: a two-chain
would follow with five source-containing sets whose isolate traces form
a maximal chain beginning at the empty trace. Let \texttt{b\_j} count
members of \texttt{B} containing exactly \texttt{j} of the three high
isolates. The sum of their frequencies gives
\texttt{b\_1+2b\_2+3b\_3\ \textless{}=\ 3(b-1)/2}, which rearranges to
\texttt{b\_1\textgreater{}=b\_2+3b\_3+3-3b\_0\textgreater{}=3(L+1)}, as
\texttt{b\_0\textless{}=1} and the full source-avoiding top and the
\texttt{L} co-leaf tops contribute to \texttt{b\_3}. These constraints
do not complete the remaining case.

\section{Verification and
dependencies}\label{verification-and-dependencies}

Run \texttt{python3 traces.py} for the finite trace check. The
multiple-source argument uses the checked pair-design, link-closure,
compressed Sauer and extra-witness lemmas from the earlier development.
The atom bound is from the written three-isolate proof. The one-source
reduction uses the checked containing-fiber, co-singleton and height-two
bounds. The companion Markdown note lists the precise Lean declarations.
The new mathematical composition is not checked end to end in Lean.

Small ground sizes in the multiple-source proof use
\href{https://ipsitransactions.org/journals/papers/tir/2017jan/fullPaper.pdf}{Vučković
and Živković, \emph{The 12-Element Case of Frankl\textquotesingle s
Conjecture}}.
\paragraph{Contributions.}
Aafi Khan Malek selected and directed the research. ChatGPT helped
develop the proofs, verify the computation, and prepare the note.
Codex was not used. This draft has not been peer reviewed.
\end{document}
